\section{Coordination and collusion}\label{sec:collusion}

An individual weight cap does not stop a cartel of many modest members. The design therefore
discounts correlated reviewers as a group, in the spirit of quadratic funding~\cite{buterin2019flexible}.

\subsection{The sublinear discount}

Reviewers are clustered by the connected components of the graph that joins $u$ and $v$ when the
correlation of their judgment vectors satisfies $|\rho_{uv}|\ge\rho_0$. A cluster $G$ with raw
weight $s_G=\sum_{u\in G}w_u$ contributes $\min\bigl(s_G,\,s_G^{\alpha}\bigr)$, with $\alpha=0.5$,
split among members in proportion to their raw weight. Each member's weight is therefore
multiplied by $\min(1,s_G^{\alpha-1})$. The engine implements and tests this discount; the
protocol does not yet apply it to the bridging weights.

\begin{lemma}\label{lem:discount}\proved{}
The discount never increases a weight. A cluster of $k$ members of unit weight counts $k^{\alpha}$
(500 coordinated reviewers count as $\sqrt{500}\approx22$), and a singleton of weight at most $1$
is unchanged.
\end{lemma}
\begin{proof}
$\min(1,s^{\alpha-1})\le1$. For $s=k\ge1$ the factor is $k^{\alpha-1}$, and $k\cdot k^{\alpha-1}=k^\alpha$.
For a singleton with $w\le1$, $s^{\alpha-1}\ge1$, so the factor is $1$.
\end{proof}

The $\min$ matters. Without it a cluster with total weight below $1$, which is every honest
singleton with $E_u\in(0,1)$, would be \emph{boosted}; the repository audit found and fixed this
defect.

\subsection{The discount is only as strong as its detector}

\begin{proposition}[Evasion pays]\label{prop:split}\proved{}
Let a cartel of total weight $s\ge1$ be detected as $j$ clusters of equal weight $s/j\ge1$. Its
total influence $j\,(s/j)^{\alpha}=j^{1-\alpha}s^{\alpha}$ is strictly increasing in $j$ for
$\alpha<1$. It reaches the undiscounted $s$ when every member escapes detection.
\end{proposition}

The discount thus converts detection recall directly into protection, and every evasion of the
detector pays. The repository audit reports that small independent noise ($\sigma=0.05$) added
to identical votes already defeats the correlation clustering. Two further weaknesses are known.
Honest like-minded reviewers are discounted as if they were a cartel. Connected components chain,
so intermediaries can merge unrelated groups, which invites griefing. A robust coordination
statistic, with a false-positive rate stated for honest like-minded reviewers, is open \openq.

\subsection{Within an epoch there is almost nothing to correlate}

\begin{proposition}[Sparse overlap]\label{prop:overlap}\proved{}
If each reviewer rates $r$ of the $m$ items of an epoch, chosen uniformly at random, then the
number of items two given reviewers share is hypergeometric with mean $r^2/m$.
\end{proposition}

With $r=9$ the expected overlap is $1.62$, $0.81$ and $0.27$ items per epoch for $m=50$, $100$
and $300$. Accumulating $30$ shared items, a modest basis for a correlation, takes about
$30m/r^2$ epochs: $19$, $37$ and $111$ respectively. The correlation matrix is therefore
undefined within an epoch at design scale. The discount can act only on long histories, while a
coalition can form and act within a single epoch. Per epoch, the operative defence is random
assignment (\cref{ass:assign}). \Cref{tab:assignment} gives the probability that a coalition
holding a share $s$ of the eligible reviewers wins a majority of a panel of nine drawn uniformly.
Stratified assignment further confines a one-camp coalition to its camp's share of the seats.
Under bridging, a one-camp majority is in any case not sufficient, although \cref{prop:leak}
shows that it helps. The threat that matters is a coalition that spans the divide, which is what
the capture-cost experiment of \cref{sec:levela} measures.

\begin{table}[t]
\centering\small
\caption{Probability that a coalition holding a share $s$ of the eligible reviewers holds at
least 5 of the 9 seats of a uniformly drawn panel. Source: \texttt{scripts/assignment.py}.}
\label{tab:assignment}
\begin{tabular}{@{}lccccc@{}}
\toprule
share $s$ & 0.1 & 0.2 & 0.3 & 0.4 & 0.5\\
\midrule
$P(\ge5\text{ of }9)$ & 0.0009 & 0.020 & 0.099 & 0.267 & 0.500\\
\bottomrule
\end{tabular}
\end{table}
